\section{Committing clearings and restoring the stage bounds}
\label{sec:assembly}

The preceding section prepares a fixed core under explicit geometry and
conditional residual laws. We now specify the scheduled clearings, verify
those inputs, and prove that every successful operation preserves the
bounds needed by later trials. The removed terminal groups are separated,
their inverse changes have controlled spatial overlap, and the changes can
be compared in one energy form fixed at the start of the stage. A failed trial is rolled back
and contributes no analytic error to this assembly.
Fix the dimension $d\ge3$ throughout.  Recall that
\[
 \nu=10^{-5}\delta/d,\qquad d\nu=10^{-5}\delta.
\]
Constants in this section may depend on $d$, but not on the final
resolution, energy in $I_d$, or a retained history reaching a committed state.

Here is the schedule for these trials. After the ordinary sweep from
stage $j-1$ to stage $j$, run clearings only when
$j\le\lfloor m/2\rfloor$. Freeze the graph whose vertices are the
pending active sites and whose edges join pairs at distance at most
$R_j=4^{8j}$. We call its connected components the snapshot components.
Fix a total order on the torus sites; the anchor of a nonempty component
is its least site. In anchor order, attempt each component
$B$ with $|B|\le t_j^{-\nu}$ once, and leave larger components alone.
If its preparation reaches DONE, reveal and record any still-unknown true
values on $B$, apply the deterministic core completion, and commit the
core together with the heat pivots as one successful clearing. This step
uses the standing invertibility of true principal matrices and adds no further
test. If preparation reaches FAIL, restore the pretrial matrix and partial
potential data, undoing all provisional eliminations. Retain every queried
answer, then mark the core and all inspected exterior sites pending, including the
pivot whose heat first failed. After all these attempts, reset the
diagonal before the next ordinary sweep. If no clearing is scheduled,
reset immediately after the ordinary sweep. The final residual sweep
has no reset. In this section the post-sweep stage is the new index
$j$, and its parameter is $t=t_j$.

Throughout a clearing stage, $t$, $r$, $p=\Im r$, and
$b=t_{n_0}t$ are frozen. We distinguish a committed prestate from a
provisional state inside a trial. The former satisfies the comparison
bounds for $D$ and the norm bound for $K$. At a running heat prestate we use
only positive bounded $p$, the stopped lag and fourth-row bounds of
Proposition~\ref{prop:clearing-trial}, the Ward identity, and the geometry
proved below. In particular, we do not assume a comparison for $D$ or a
norm bound for $K$ at a provisional heat prefix. The committed bounds are
proved below from the outputs of successful operations and the stage hypotheses; they are not additional
stopping tests for the preparation.

\subsection{Geometry and localization during a clearing}

For a terminal set, give a nearest-neighbor edge cost zero if either
endpoint is terminal, and cost one otherwise. Its path pseudometric is
denoted by $d_T$; $\dist$ always denotes ordinary torus distance.
In particular, the subscript $T$ distinguishes the collapsed metric
from the dimension $d$.

\begin{lemma}[Adding separated terminal groups]
\label{lem:assembly-metric}
Let $c>0$ and $\mu\ge0$, and suppose
$d_T(u,v)\ge c\dist(u,v)-\mu$. Let $O_a$ be sets of new terminals,
each containing at most $M$ sites, with
$\dist(O_a,O_{a'})\ge R-2\ell$ for $a\ne a'$. If
\[
 \mathcal L=c(R-2\ell)-\mu>2M,
\]
then the pseudometric $d_T'$ after all these additions satisfies
\[
 d_T'(u,v)\ge c(1-2M/\mathcal L)\dist(u,v)-\mu-2M.
\]
The number of added groups does not enter this bound.
\end{lemma}

\begin{proof}
A path minimizing the new cost can be chosen simple by deleting cycles.
Suppose it meets $k\ge1$ different groups, and let $A$ be its old cost.
Each site of their union is incident to at most two edges of the simple
path. Its cost therefore decreases by at most $2Mk$, even if the path
returns to a group after leaving it. In the chronological sequence of
group visits there are at least $k-1$ transitions between different
groups. To obtain disjoint segments, list the groups visited along the
path, suppress consecutive repetitions of the same group, and use the
segment from the last visit to one group before the first visit to the
next. There are at least $k-1$ such transitions. Each has old cost at
least the old collapsed distance between its endpoints, hence at least
$\mathcal L$. This argument already includes all shortcuts through
old terminals. Repeated visits can increase the number of transitions
but do not increase the count of distinct groups in the loss $2Mk$.
Hence $A\ge(k-1)\mathcal L$ and
\[
 d_T'(u,v)\ge A-2Mk
 \ge (1-2M/\mathcal L)A-2M
 \ge c(1-2M/\mathcal L)\dist(u,v)-\mu-2M.
\]
If the path meets no new group, the old lower bound applies directly.
\end{proof}

\begin{lemma}[Localization of a clearing trial]
\label{lem:clearing-locality}
At post-sweep stage $j\le\lfloor m/2\rfloor$, put
\[
 R_j=4^{8j},\qquad \ell_j=4^{4j},\qquad
 r_*=t_j^{-\nu},\qquad n_h=\lfloor t_j^{-1-e_h}\rfloor,
 \qquad e_h=.46\delta.
\]
Take the snapshot components of pending sites joined at distance at most
$R_j$, and consider only components $B$ of size at most $r_*$.
Assume the terminal set was initially empty and the preceding clearing
stages used the same cardinality and pivot-location rules. The new terminals
of a completed successful group consist of its core and heat pivots. The
current provisional group contains its heat pivots, with the next proposed
pivot included when its prospective output is considered. Suppose each
such counted group has at most $Cn_h$ sites and all previous pivots in a
trial lie within $\ell_j/2$ of its snapshot component. The following
assertions hold after increasing $n_0$ uniformly in $m$:
\begin{enumerate}
\item At every committed endpoint and running heat prestate, the
terminal metric satisfies
\begin{equation}
 d_T(u,v)\ge\tfrac12\dist(u,v)-t_j^{-1.03}.
 \label{eq:assembly-terminal-metric}
\end{equation}
\item At a running heat prestate, every proposed heat pivot, including a pivot whose update
would first break a distortion budget, lies within $\ell_j/2$ of $B$.
After that pivot is included as a provisional terminal, the metric bound
\eqref{eq:assembly-terminal-metric} still holds at its prospective output,
even if the output first breaks a cap or reaches the heat limit.
\item At each committed endpoint, running heat prestate, and prospective output,
at least $N/2$ physical
sites remain active. In particular, the current exterior of an eligible
core is nonempty for sufficiently large $n_0$.
An eligible component and all its edits lie within distance
$D_s=Ct_j^{-\nu}R_j$ of its anchor. Cubes of radii $100D_s$ or $200D_s$
about these anchors have overlap at most $Ct_j^{-d\nu}$.
\end{enumerate}
These conclusions require no bound on $K$ or comparison for $D$ at
provisional heat prefixes.
\end{lemma}

\begin{proof}
We first prove the metric assertion under the stated location hypothesis,
and then use it to continue that hypothesis to the next pivot.
Distinct snapshot components are separated by more than $R_j$.
Their new terminal sets, including the current provisional additions,
are therefore separated by at least $R_j-2\ell_j$. For an actual success,
the number of new terminals is less than $n_h+|B|\le n_h+r_*\le Cn_h$
for small $t_j$, since $\nu<1+e_h$. A provisional group, including its
one prospective pivot, obeys the same bound from the heat limit.
Apply Lemma~\ref{lem:assembly-metric} with $M=Cn_h$ to all additions of
the stage at once. Starting with the ordinary metric before any clearing,
the multiplicative losses are bounded by a constant times
\[
 \sum_{j\ge n_0}\frac{t_j^{-1-e_h}}{R_j}
 \le C\sum_{j\ge n_0}4^{-j(7-e_h)}.
\]
This sum can be made arbitrarily small. The additive defects, through
stage $j$, are at most
\[
 C\sum_{i=n_0}^j t_i^{-1-e_h}
 \le C't_j^{-1-e_h}\le t_j^{-1.03},
\]
where $e_h=.0046<.03$. In these estimates $t_i\asymp4^{-i}$ uniformly
in the final resolution. They also imply the inequality $\mathcal L>2M$
required in the metric lemma. This proves
\eqref{eq:assembly-terminal-metric}.

At the current running heat prestate, the metric just proved supplies
\eqref{eq:clearing-distance-input} with $\ell=\ell_j\asymp t_j^{-4}$.
The selection implication \eqref{eq:clearing-proposed-location} therefore
places the next proposed pivot within $\ell_j/2$ of $B$. Its underlying
full-range estimate uses only survivor damping and terminal Schur structure,
so this application needs neither a lower singular-value bound for the
cross block nor a provisional bound on $K$ or $D$.
Induction over the proposed pivots proves that there is no first
violation of the location hypothesis. In particular the argument
applies before a prospective first budget-breaking update; it does not
assume that update will be accepted.
Once that pivot is localized, applying the same separated-group bound
with it included also gives the metric estimate at its prospective
poststate. A subsequent rollback removes its numerical metric cost.

One anchor from each snapshot component forms an $R_j$-separated set.
Packing disjoint balls of radius a fixed fraction of $R_j$ bounds the
number of components by $CN/R_j^d$. Each success, or the current
provisional trial, removes at most $Cn_h$ sites. The total fraction
removed through all stages is bounded by
\[
 C\sum_{j\ge n_0}n_h/R_j^d
 \le C\sum_{j\ge n_0}t_j^{8d-1-e_h}<\tfrac12
\]
for large $n_0$. Failed trials are rolled back and have no terminal
cost. This proves the active-count assertion, including provisional
prefixes and their prospective outputs. Moreover,
$r_*=t_j^{-\nu}\le C4^{\nu m/2}\ll N=4^{20dm}$ for $j\le m/2$.
Removing the at most $r_*$ core sites from an active set of size at least
$N/2$ leaves a nonempty exterior.
At DONE, the precompletion terminal and active sets are those of the
committed state at trial start if no heat occurred, and those of the last
heat output otherwise. The same metric and active-count bounds therefore
hold there.

A snapshot component of size at most $r_*$ has diameter at most
$(r_*-1)R_j$. Adding the pivot neighborhood gives the stated $D_s$.
Assign exactly one cube to each snapshot component, centered at its
anchor. For all cubes containing a given point, the disjoint balls of
radius $R_j/3$ about their anchors fit in a ball of radius $CD_s$.
The $d$-dimensional volume bounds therefore give overlap at most
$C_d(D_s/R_j)^d=C_dt_j^{-d\nu}$. There is no factor for the number of
sites in a component: all its edits use this single enclosing cube.
All these cubes fit on the torus: $j\le m/2$ implies
$D_s\le C4^{(8+\nu)m/2}\ll4^{20m}$.
\end{proof}

\paragraph{The true residual at a selected pivot.}
The geometry now identifies which true information is available in an
actual scheduled trial. All ordinary pending marks are made before the
snapshot is frozen. A pending snapshot site outside $B$ therefore belongs
to a different component, at distance greater than $R_j$ from $B$.
An exterior inspected in an earlier failed clearing lies within
$\ell_j/2$ of its own snapshot component, hence at distance greater than
$R_j-\ell_j/2$ from $B$. Neither kind of site can be the current pivot,
which lies within $\ell_j/2$ of $B$. Earlier successful inspections are
committed terminals and unavailable; earlier pivots in this trial have
been provisionally removed. The ordinary-sweep rule has preserved the
invariant that a coupled active site has an untouched true suffix. The
exclusions above therefore show that the selected localized exterior pivot is
coupled and has an untouched true residual. This conclusion includes the pivot
whose update first causes failure.

The pivot and ridge are measurable from the retained pre-query history:
the current provisional matrix is known, and the ridge and tie rules use
measurable spectral functions of that matrix. Lemma~\ref{lem:conditional-freshness}
therefore gives the conditional law of the untouched residual $\xi$ at
the selected site. At stage $j$ its moments are exactly
\[
 \EE(\xi\mid\mathcal F)=0,\qquad
 \EE(\xi^2\mid\mathcal F)=s_j=t_j-\eta,\qquad
 |\xi|\le\sqrt{3s_j},\qquad \eta/t_j\le Ct_j,
\]
where the last inequality uses $j\le m/2$. The deterministic location
argument did not assume this product law; it first proved that the chosen
coordinate is among those to which the law applies.

Every exposed residual is recorded before its outcome is tested. On FAIL,
rollback restores the matrix and partial potential data from trial start
before marking the core and every inspected exterior, including the first
breaking pivot, while retaining their answers. On DONE, a pending core may still contain unknown true
values. Revealing and recording those values makes the completed numerical
state a function of the retained history. Lemma~\ref{lem:core-completion}
applies pointwise to the obtained values under its standing invertibility
hypothesis; it needs no centered law or freshness for them. Core and heat
pivots then become committed terminals in the one successful operation.
Thus every active site queried during this clearing becomes terminal on success or pending on
rollback; no exposed suffix remains at a coupled active site for the next
component or stage.
The invertibility of true principal matrices is the probability-one fact
at the fixed energy from Section~\ref{sec:state}, not a preliminary query or an additional
conditioning of the product law.

These observations verify the geometric and random-input hypotheses of
Proposition~\ref{prop:clearing-trial} for each scheduled eligible component.
Once the deterministic assembly below supplies its trial-start lag and row
slack at every committed prefix, the conditional failure estimate applies
after any earlier failed trials, with all their information retained.

\subsection{Summing the committed changes}

A successful clearing supplies one difference between its inverse before
and after the operation, restricted to its surviving sites. We first state the deterministic result that sums
these differences. The energy error has two parts: the change of the
surviving edges, and the old energy carried through the sites that have
been removed. Both will be compared in the spatial form fixed at the start
of the stage.

\begin{proposition}[Deterministic assembly of committed changes]
\label{prop:committed-assembly}
Fix a clearing stage with parameter $t$. Consider the finite sequence of
its committed numerical states: the initial post-sweep state, indexed by
$0$, and the states after successive successful clearings, indexed by
$a=1,\ldots,k_{\max}$. Failed trials leave this numerical sequence
unchanged. Let $A_a\subset A_{a-1}$ be the active sites after operation
$a$, put $O_a=A_{a-1}\setminus A_a$, and let $T_a$ be the complement of
$A_a$ in the physical torus.

At each index use the physical state of Section~\ref{sec:state}: the
sites in $T_a$ have their exact true potentials, the sites in $A_a$ have
their current partial potentials, and $M^{(a)}$ is the real terminal-Schur
matrix. The relevant true terminal principal matrices are invertible at
the fixed energy. Use the restriction of one frozen vector $r$ at every
index, with $p=\Im r$ bounded above and below by positive constants, and
write
\[
 g^{(a)}=(M^{(a)}-t\diag(r|_{A_a}))^{-1},\qquad
 h^{(a)}=\sqrt t\,g^{(a)},\qquad
 K^{(a)}_{ij}=(h^{(a)}_{ij})^2,\qquad
 F^{(a)}_{ij}=|h^{(a)}_{ij}|^2.
\]
Let $D_a$ be the Ward Laplacian of this state, so
\[
 D_a(f)=\frac12\sum_{i,j\in A_a}p_ip_jF^{(a)}_{ij}|f_i-f_j|^2.
\]
With $q_i^{(a)}=\Im g^{(a)}_{ii}$, the Ward identity gives
$F^{(a)}p=q^{(a)}$.
The word committed here identifies the numerical endpoints. The comparison
at the initial endpoint is an input; its transport to later endpoints is
the conclusion below.

Assume at the initial state that, on potentials modulo constants,
\begin{equation}
 cW\le D_0\le CW,\qquad
 \max_{i\in A_0}\sum_{j\in A_0}|h^{(0)}_{ij}|^4
 \le .7H_*,\qquad
 \max_{i\in A_0}|g^{(0)}_{ii}-r_i|\le Ct^{.42\delta}.
 \label{eq:assembly-initial-data}
\end{equation}
Here $H_*=t^{1+\delta}$ and
$W(f)=\inf_{f=\alpha+v|_{A_0}}
 (\|\alpha\|_2^2+b^{-1}\|\nabla v\|_2^2)$, with $b=t_{n_0}t$.

Assume also the geometric bounds supplied by
Lemma~\ref{lem:clearing-locality}. For $D_s\asymp t^{-8-\nu}$, each
$O_a$ lies in the ball of radius $D_s$ about its anchor $z_a$ and has
size at most $Cn_h$, where $n_h=\lfloor t^{-1-.46\delta}\rfloor$.
The anchors are distinct. The cube halos $Q_a,Q'_a$
of radii $100D_s,200D_s$ are embedded lattice cubes in the torus, and
each family has overlap at most $Ct^{-d\nu}$. At every indexed state,
$|A_a|\ge N/2$ and
\[
 d_{T_a}(u,v)\ge\tfrac12\dist(u,v)-m',\qquad m'\le t^{-1.03}.
\]

For each operation set
\[
 a^{(a)}=h^{(a)}-h^{(a-1)}|_{A_a\times A_a}.
\]
Assume these actual survivor differences satisfy
\begin{align}
 \max_{i,j\in A_a}|a^{(a)}_{ij}|
 &\le C\sqrt t\,t^{.42\delta},\notag\\
 \max_{i\in A_a}\sum_{j\in A_a}|a^{(a)}_{ij}|^2
 &\le Ct^{.25\delta},\qquad
 \sum_{i,j\in A_a}|a^{(a)}_{ij}|^2
 \le Ct^{-1-.47\delta}.
 \label{eq:assembly-single-distortions}
\end{align}
Assume further that, with constants $c,C>0$ and a fixed $P$,
\begin{equation}
 |a^{(a)}_{ij}|
 \le t^{-P}\exp\!\left\{-ct\left[
 (\dist(i,O_a)-Cm')_++(\dist(j,O_a)-Cm')_+
 \right]\right\}.
 \label{eq:assembly-single-decay}
\end{equation}
Then, for all sufficiently small $t$, every prefix $0\le k\le k_{\max}$
satisfies
\begin{align}
 \|K^{(k)}\|&\le\|K^{(0)}\|+t^{.10\delta},
 \label{eq:assembly-K-output}\\
 \max_{i\in A_k}|g^{(k)}_{ii}-r_i|
 &\le Ct^{.42\delta-d\nu}=o(t^{.36\delta}),
 \label{eq:assembly-lag-output}\\
 \max_{i\in A_k}\left(\sum_{j\in A_k}|h^{(k)}_{ij}|^4\right)^{1/4}
 &\le\left(.7^{1/4}+Ct^{.0225\delta-d\nu}\right)H_*^{1/4}
 <(.8H_*)^{1/4}.
 \label{eq:assembly-row-output}
\end{align}
All Schur complements in the successive definition
\[
 L_a=D_a-\Schur_{O_a}D_{a-1}\quad\hbox{on }A_a
\]
exist, and each $D_a$ has precisely the constant kernel. Extend each
$L_a$ by zero to $A_0$. Then
\begin{equation}
 \left|\sum_{a=1}^k L_a\right|
 \le t^{.10\delta}D_0
 \label{eq:assembly-energy-output}
\end{equation}
in quadratic forms.  In particular, writing $X_k=D_k^+$, every neutral
charge $y$ supported in $A_k$ satisfies
\begin{equation}
 (1+t^{.10\delta})^{-1}X_0(y)
 \le X_k(y)\le
 (1-t^{.10\delta})^{-1}X_0(y),
 \label{eq:assembly-inverse-output}
\end{equation}
where $y$ is extended by zero in $X_0(y)$.  All constants are uniform
in the number of components and the torus volume.
\end{proposition}

\paragraph{The inputs supplied by a successful clearing.}
We verify the two per-operation bounds before proving the deterministic
summation. At a trial satisfying the starting hypotheses of
Proposition~\ref{prop:clearing-trial}, its DONE output gives the heat
budgets and the final projection cap. By the schedule above, any
still-unknown core values are then revealed and recorded, and
Lemma~\ref{lem:core-completion} applies pointwise; the nonempty-exterior
condition is supplied by Lemma~\ref{lem:clearing-locality}.
On the final survivors, the recorded heat difference is exactly
$\hat h-h^0$ restricted to those survivors. Adding the core difference
gives the actual endpoint difference $a^{(a)}$.

The heat contribution to
\eqref{eq:assembly-single-distortions} has entry bound
$\sqrt t\,t^{.42\delta}$, squared row bound $t^{.25\delta}$,
and squared Hilbert--Schmidt bound
$n_h t^{-.01\delta}\le t^{-1-.47\delta}$.
The corresponding core bounds are
$C\sqrt t\,t^{.44\delta}$, $Ct^{.44\delta}$, and
$Cr_*/t=Ct^{-1-\nu}$.  The triangle inequality, squared when
necessary, gives the displayed combined bounds because
$\nu<.47\delta$. Restricting the heat record to the final survivors can
only decrease its entry, row, and Frobenius bounds; the core difference
is already indexed by the final survivors.

To verify \eqref{eq:assembly-single-decay}, consider a pair of final
survivors. A heat at $x\in O_a$ changes its entry by
$c_xh_{ix}h_{xj}$ with $|c_x|\le C$.
Lemma~\ref{lem:locality} and \eqref{eq:assembly-terminal-metric}
bound this product by
\[
 Ct^{-1}\exp\{-ct[(\dist(i,O_a)-Cm')_+
                         +(\dist(j,O_a)-Cm')_+]\}.
\]
There are at most $n_h$ heats, so their sum has the claimed form.
For the core change, Corollary~\ref{cor:core-decay} gives the same decay
from the core, and the core is contained in $O_a$; use the same terminal
metric.
The prefactors, including $n_h$, are bounded by a fixed power of $t^{-1}$.
This proves the asserted decay without any comparison for $D$ at a
provisional prefix. Thus each successful scheduled operation supplies
the two bounds assumed in the proposition, provided its trial starts with
the stated lag and row slack. The prefix conclusions below supply that
slack to the next trial when the result is used in the stage induction.

\begin{proof}[Proof of Proposition~\ref{prop:committed-assembly}]
Put $\epsilon=t^{.10\delta}<1/2$. We prove the form bound and the
existence of its Schur complements by induction on the prefix length $k$.
In the step for $k$, suppose the form bound and the identity below hold
at every shorter prefix. For each $a\le k$, let
$S_{a-1}=\sum_{s<a}L_s$, padded to $A_0$, and
$E_{a-1}=A_0\setminus A_{a-1}$. The induction hypothesis is
$|S_{a-1}|\le\epsilon D_0$ together with
\begin{equation}
 D_{a-1}=\Schur_{E_{a-1}}(D_0+S_{a-1}).
 \label{eq:assembly-schur-identity}
\end{equation}
For the empty prefix, $S_0=0$ and $E_0=\varnothing$, so the identity is immediate.
The initial lower comparison with $W$ makes the kernel of $D_0$
precisely the constants. At a later prefix, the form bound gives the same
kernel for $D_0+S_{a-1}$. The Schur fact from Section~\ref{sec:state}
then gives the constant kernel of $D_{a-1}$ and, since $A_a$ is nonempty,
the existence of its pivot on $O_a$. Thus $L_a$ is defined. Minimization
over eliminated coordinates also gives, for every extension $F$ to $A_0$,
\begin{equation}
 D_{a-1}(F|_{A_{a-1}})
 \le(1+\epsilon)D_0(F)\le CW(F).
 \label{eq:assembly-pre-comparison}
\end{equation}

\paragraph{The two energy contributions.}
On surviving pairs the changes of $K$ and $F$ are exactly
\[
 \Delta K^{(a)}=2h^{(a-1)}\circ a^{(a)}+a^{(a)}\circ a^{(a)},
 \qquad
 \Delta F^{(a)}=2\Re(\overline{h^{(a-1)}}\circ a^{(a)})
                    +|a^{(a)}|^2,
\]
where old matrices are restricted to $A_a\times A_a$ and $\circ$ denotes
entrywise multiplication. Let $D^{\rm edge}_{A_a}$ be the Laplacian made
from the old edges with both endpoints in $A_a$. For a potential $f$ on
$A_a$, the fact that $p$ is frozen gives
\begin{equation}\label{eq:assembly-energy-split}
\begin{aligned}
 (D_a-D^{\rm edge}_{A_a})(f)
 &=\frac12\sum_{i,j\in A_a}p_ip_j\Delta F^{(a)}_{ij}|f_i-f_j|^2,\\
 B_a&=\Schur_{O_a}D_{a-1}-D^{\rm edge}_{A_a}\ge0,\\
 L_a&=(D_a-D^{\rm edge}_{A_a})-B_a.
\end{aligned}
\end{equation}
The first term uses the single endpoint difference $a^{(a)}$. The form
$B_a$ is nonnegative because the Schur complement minimizes the old
energy over $O_a$, while the surviving-edge energy does not depend on
those coordinates. It carries old energy through the removed sites, so
its far part will require the old resolvent decay separately from the
decay of $a^{(a)}$.

Fix a real potential $f$ on $A_0$ and one split
$f=\alpha+v|_{A_0}$ for the whole prefix. Put
\[
 c_a=\operatorname{avg}_{Q'_a}v,\qquad
 W_a=\|\alpha\|_{2,Q'_a\cap A_0}^2
       +b^{-1}\|\nabla v\|_{2,Q'_a}^2.
\]
The last norm sums nearest-neighbor edges of the cube; harmless enlargement
of $Q'_a$ includes its boundary edges. We use this same split for every
$a$ and take the infimum defining $W(f)$ only after summing the local
estimates.

\paragraph{Uniform tails for the survivor changes.}
We first record the tail estimate for $a^{(a)}$ and the two survivor
changes it produces. Put $\rho_i=\dist(i,z_a)$ and
$\sigma_i=(\dist(i,O_a)-Cm')_+$.  Since $O_a$ lies within distance
$D_s$ of $z_a$ and $Cm'=o(D_s)$,
\[
 \sigma_i+\sigma_j\ge \rho_i+\rho_j-4D_s.
\]
If $(i,j)\notin Q_a\times Q_a$, then $\rho_i+\rho_j\ge100D_s$,
and hence
\[
 \sigma_i+\sigma_j\ge\tfrac12(\rho_i+\rho_j)+46D_s.
\]
It follows from \eqref{eq:assembly-single-decay}, after decreasing
$c$ and increasing $P$, that
\begin{equation}
 |a^{(a)}_{ij}|\mathbf1_{(i,j)\notin Q_a^2}
 \le t^{-P}e^{-ctD_s}e^{-ct(\rho_i+\rho_j)}.
 \label{eq:assembly-anchor-tail}
\end{equation}
The inverse of a real symmetric active matrix with damping
$-it\diag(p)$ has norm at most $C/t$.  Thus
$|h^{(a-1)}_{ij}|\le Ct^{-1/2}$, and
\eqref{eq:assembly-anchor-tail} also holds for the absolute values
of the tail entries of $\Delta K^{(a)}$ and $\Delta F^{(a)}$, after
increasing $P$.

On the $d$-dimensional torus, uniformly in its side length, for every
fixed $q\ge0$ and $0<t\le1$,
\begin{equation}
 \sum_z e^{-ct\dist(i,z)}\dist(i,z)^q
 \le C_{d,q}t^{-d-q}.
 \label{eq:assembly-lattice-moments}
\end{equation}
For $q=0$ take $0^0=1$. Indeed, counting coordinate representatives
of minimal length gives at most $C_d(1+s)^{d-1}$ sites in the shell
$s\le\dist(i,z)<s+1$. The sum is bounded by
$C_{d,q}\sum_{s\ge0}e^{-cts}(1+s)^{d+q-1}$, which is at most the
right side by comparison with its integral. In particular the zeroth
and second moments cost $Ct^{-d}$ and $Ct^{-d-2}$.
Distinctness of the anchors and
the triangle inequality consequently give
\[
 \sum_a e^{-ct(\dist(i,z_a)+\dist(j,z_a))}
 \le Ct^{-d}e^{-ct\dist(i,j)/2}.
\]
Therefore the sum of all tail kernels has row and column sums at
most $t^{-P}e^{-ctD_s}$.  The same conclusion holds after arbitrary
restrictions to survivors or to a committed prefix.

To use the same tail in the energy form, extend the fixed $\alpha$ by
zero to the torus. For every
lattice displacement $u$, a shortest nearest-neighbor path and
Cauchy--Schwarz give
\begin{equation}
 \sum_i|v_{i+u}-v_i|^2
 \le\dist(0,u)^2\|\nabla v\|_2^2.
 \label{eq:assembly-translation}
\end{equation}
Apply the row-sum bound to the $\alpha$ differences and
\eqref{eq:assembly-translation} to the $v$ differences.  Equations
\eqref{eq:assembly-anchor-tail}--\eqref{eq:assembly-lattice-moments}
then imply
\begin{equation}
 \sum_a\sum_{(i,j)\notin Q_a^2}
 |\Delta F^{(a)}_{ij}|\,|f_i-f_j|^2
 \le t^{-P}e^{-ctD_s}
       (\|\alpha\|_2^2+\|\nabla v\|_2^2).
 \label{eq:assembly-tail-form}
\end{equation}
We henceforth denote such bounds by $\varepsilon_t$ times the
indicated norm, where $\varepsilon_t=t^{-P}e^{-ctD_s}$ with $P$
allowed to increase.  Because $tD_s\asymp t^{-7-\nu}$,
$\varepsilon_t$ is smaller than every fixed positive power of $t$.
Since $b\le1$, the right side of \eqref{eq:assembly-tail-form} is
at most $\varepsilon_t(\|\alpha\|_2^2+b^{-1}\|\nabla v\|_2^2)$.

\paragraph{The diagonal, fourth powers, and $K$.}
For a site surviving a prefix, its diagonal increments telescope.
There are at most $Ct^{-d\nu}$ halos $Q_a$ containing that site.
Their near increments obey
$|\Delta g^{(a)}_{ii}|\le Ct^{.42\delta}$; all other increments
sum to at most $\varepsilon_t$, by
\eqref{eq:assembly-anchor-tail} and $\Delta g=a^{(a)}/\sqrt t$.
Together with \eqref{eq:assembly-initial-data}, this proves
\eqref{eq:assembly-lag-output}.  In particular, at every committed
prestate the Ward identity gives
\begin{equation}
 \sum_j|h_{ij}|^2\le q_i/\min_jp_j\le C.
 \label{eq:assembly-row-two}
\end{equation}

The first two bounds in \eqref{eq:assembly-single-distortions} imply
\[
 \sum_j|a^{(a)}_{ij}|^4
 \le \max_j|a^{(a)}_{ij}|^2\sum_j|a^{(a)}_{ij}|^2
 \le Ct^{1+1.09\delta}.
\]
For the near pieces, Minkowski's inequality in $\ell^4$, followed
by halo overlap, costs at most
\[
 Ct^{-d\nu}t^{(1+1.09\delta)/4}
 =Ct^{.0225\delta-d\nu}H_*^{1/4}
\]
per row.  The sum of the $\ell^4$ norms of the tail rows is at most
their summed $\ell^1$ norm, which is bounded by $\varepsilon_t$.
Telescoping on the final surviving columns, and observing that
discarding columns decreases the initial row norm, proves
\eqref{eq:assembly-row-output}.  The strict final inequality holds
because $.0225\delta-d\nu>0$.

For $K$, at a committed prestate
\eqref{eq:assembly-row-two} and Cauchy--Schwarz yield
\[
 \sum_j|\Delta K^{(a)}_{ij}|
 \le2\left(\sum_j|h_{ij}|^2\right)^{1/2}
         \left(\sum_j|a^{(a)}_{ij}|^2\right)^{1/2}
       +\sum_j|a^{(a)}_{ij}|^2
 \le Ct^{.125\delta}.
\]
The matrices are symmetric, so the same bound holds for column
sums.  The near pieces have at most $Ct^{-d\nu}$ overlapping
row supports, and the tails have norm at most $\varepsilon_t$.
On $A_k$ the exact telescoping identity is
\[
 K^{(k)}=K^{(0)}|_{A_k\times A_k}
       +\sum_{a\le k}\Delta K^{(a)}|_{A_k\times A_k}.
\]
The row and column sum bound therefore gives
$\|K^{(k)}\|\le\|K^{(0)}\|+Ct^{.125\delta-d\nu}
+\varepsilon_t$, which implies \eqref{eq:assembly-K-output}.

\paragraph{The local energy of a survivor change.}
We now estimate the near part of the first term in
\eqref{eq:assembly-energy-split} using the common split fixed above.
For nonnegative symmetric weights $w_{ij}$ supported in $Q_a^2$,
let $l_i=\sum_jw_{ij}$, $L_\infty=\max_i l_i$ and
$L_1=\sum_i l_i$.  The difference inequality and
Lemma~\ref{lem:cube-weight} give
\begin{align}
 \sum_{i,j}w_{ij}|f_i-f_j|^2
 &\le C L_\infty\|\alpha\|_{2,Q'_a\cap A_0}^2
       +C\sum_i l_i|v_i-c_a|^2\notag\\
 &\le C\left(L_\infty+
       bL_\infty^{1/3}L_1^{2/3}\right)W_a.
 \label{eq:assembly-weighted-local}
\end{align}
Apply this with $w_{ij}=|a^{(a)}_{ij}|^2$ on surviving pairs.
By \eqref{eq:assembly-single-distortions} and $b=t_{n_0}t\le t$,
\[
 L_\infty\le Ct^{.25\delta},\qquad
 L_1\le Ct^{-1-.47\delta},\qquad
 bL_\infty^{1/3}L_1^{2/3}
 \le Ct^{1/3-.23\delta}.
\]
Since $1/3-.23\delta>.25\delta$ for $\delta=.01$, we obtain
\begin{equation}
 \sum_{i,j\in Q_a\cap A_a}
 |a^{(a)}_{ij}|^2|f_i-f_j|^2
 \le Ct^{.25\delta}W_a.
 \label{eq:assembly-delta-local}
\end{equation}

For the old factor in the cross term, apply the preceding-prefix bound
\eqref{eq:assembly-pre-comparison} with a local extension. Choose a cutoff
$\chi_a$ equal to one on $Q_a$, supported in
$Q'_a$, and with $|\nabla\chi_a|\le C/D_s$.  Set
$F=\chi_a\alpha+\chi_a(v-c_a)|_{A_0}$.  Cube Poincare gives
\[
 \|\nabla[\chi_a(v-c_a)]\|_2^2
 \le C\left(\|\nabla v\|_{2,Q'_a}^2
       +D_s^{-2}\|v-c_a\|_{2,Q'_a}^2\right)
 \le C\|\nabla v\|_{2,Q'_a}^2,
\]
and hence $W(F)\le CW_a$.  Differences of $F$ agree with
those of $f$ on $Q_a^2$.  The old Laplacian edge weights
$p_ip_j|h_{ij}|^2$ are nonnegative; therefore
\eqref{eq:assembly-pre-comparison} bounds their near commutator
energy by $CW_a$.  Weighted Cauchy--Schwarz with
\eqref{eq:assembly-delta-local} proves that the near survivor-edge
change has absolute form at most
\begin{equation}
 Ct^{.125\delta}W_a.
 \label{eq:assembly-survivor-energy}
\end{equation}
Indeed its cross term costs the square root of
$CW_a\cdot Ct^{.25\delta}W_a$, and its quadratic term costs
$Ct^{.25\delta}W_a$.  The bounded factors $p_ip_j$ only alter
the constants.

\paragraph{The old energy carried through $O_a$.}
We next estimate $B_a$, the second term in
\eqref{eq:assembly-energy-split}. Test the old Schur minimum with the
constant choice $f_x=c_a$ for all $x\in O_a$. This gives
\begin{equation}
 0\le B_a(f|_{A_a})
 \le C\sum_{i\in A_a}\sum_{x\in O_a}
 F^{(a-1)}_{ix}|f_i-c_a|^2.
 \label{eq:assembly-bridge-trial}
\end{equation}
All old edges with both endpoints in $O_a$ vanish for this trial.

For $i\in Q_a$, put $l_i=\sum_{x\in O_a}F^{(a-1)}_{ix}$.
The row estimate above gives fourth-power rows below $.8H_*$ at the old
state $a-1$. Cauchy--Schwarz and
$|O_a|\le Cn_h$ imply
\[
 \max_i l_i\le C\sqrt{n_hH_*}=Ct^{.27\delta}.
\]
By symmetry and \eqref{eq:assembly-row-two},
$\sum_i l_i\le Cn_h\le Ct^{-1-.46\delta}$.  Thus
\[
 b(\max_i l_i)^{1/3}(\sum_i l_i)^{2/3}
 \le Ct^{1/3-(13/60)\delta}
 \le Ct^{.27\delta}.
\]
Apply Lemma~\ref{lem:cube-weight} to the $v-c_a$ term in
\eqref{eq:assembly-bridge-trial}, and the supremum bound to its
$\alpha$ term.  The resulting near contribution is at most
$Ct^{.27\delta}W_a$.

The far part of this constant trial uses the old terminal geometry.
If $i\notin Q_a$ and $x\in O_a$, then
\[
 d_{T_{a-1}}(i,x)\ge\tfrac12(\dist(i,z_a)-D_s)-m'
          \ge c\dist(i,z_a).
\]
The old Combes--Thomas bound of Lemma~\ref{lem:locality} therefore
gives, after summing the at most $Cn_h$ removed sites,
\[
 \sum_{x\in O_a}F^{(a-1)}_{ix}
 \le t^{-P}e^{-ctD_s}e^{-ct\dist(i,z_a)}.
\]
In \eqref{eq:assembly-bridge-trial} use
\[
 |f_i-c_a|^2\le3\bigl(|\alpha_i|^2+|v_i-v_{z_a}|^2
                              +|v_{z_a}-c_a|^2\bigr).
\]
After summing distinct anchors, the first two terms are bounded
by \eqref{eq:assembly-lattice-moments} and
\eqref{eq:assembly-translation}, respectively.  For the last term,
cube Poincare gives the sufficient pointwise estimate
\[
 |v_{z_a}-c_a|^2
 \le\sum_{Q'_a}|v-c_a|^2
 \le CD_s^2\|\nabla v\|_{2,Q'_a}^2.
\]
Summing the far weights costs $Ct^{-d}$, and summing these local
gradient energies costs the halo overlap $Ct^{-d\nu}$.  All
remaining factors are powers of $t^{-1}$, absorbed into
$\varepsilon_t$.  Hence the total far contribution from the
Schur trials is at most
\begin{equation}
 \varepsilon_t(\|\alpha\|_2^2+b^{-1}\|\nabla v\|_2^2).
 \label{eq:assembly-bridge-tail}
\end{equation}

\paragraph{Closing the comparison induction.}
Use the decomposition \eqref{eq:assembly-energy-split}. Equations
\eqref{eq:assembly-survivor-energy},
\eqref{eq:assembly-tail-form}, and
\eqref{eq:assembly-bridge-tail}, followed by halo overlap, imply
for the current prefix and the one fixed split of $f$ that
\begin{align*}
 |S_k(f)|
 &\le Ct^{.125\delta}\sum_{a\le k}W_a
       +\varepsilon_t(\|\alpha\|_2^2+b^{-1}\|\nabla v\|_2^2)\\
 &\le\bigl(Ct^{.125\delta-d\nu}+\varepsilon_t\bigr)
       (\|\alpha\|_2^2+b^{-1}\|\nabla v\|_2^2).
\end{align*}
Take the infimum over all splits and use $W\le C D_0$.
Since
\[
 .125\delta-d\nu=.12499\delta>.10\delta,
\]
the resulting estimate is at most
$\tfrac12t^{.10\delta}D_0(f)$ for sufficiently small $t$.
This is a strict improvement of the induction bound. It also makes the
kernel of $D_0+S_k$ precisely the constants. Since $A_k$ is nonempty,
its Schur pivot on $E_k$ exists. The zero extension of $L_k$ and
associativity of Schur complements now give
\[
 \Schur_{E_k}(D_0+S_k)
 =\Schur_{O_k}(D_{k-1}+L_k)
 =\Schur_{O_k}D_{k-1}+L_k=D_k.
\]
This continues \eqref{eq:assembly-schur-identity} and gives the constant
kernel of $D_k$. Thus the form induction closes. The diagonal and row
estimates proved above also preserve the trial-start margins. The
application uses no comparison for $D$ or cap on $K$ at an intermediate
heat of a provisional trial.

Finally \eqref{eq:assembly-schur-identity} holds also at prefix
$k$.  For a neutral charge supported in $A_k$, inversion of a
Schur complement gives the inverse quadratic form of
$D_0+S_k$ on its zero extension.  The inequalities
$(1-\epsilon)D_0\le D_0+S_k\le(1+\epsilon)D_0$ therefore imply
\eqref{eq:assembly-inverse-output} by inversion on the space
modulo constants.
\end{proof}

\subsection{Resetting the diagonal}

The completed stage has changed the resolvent while keeping $r$ frozen.
We restore the relation $r_i=g_{ii}$ by a small deterministic change in
$r$. This is the only reset between consecutive ordinary sweeps.

\begin{proposition}[Diagonal reset]
\label{prop:reset}
Let $g=(M-t\diag r)^{-1}$ on a nonempty active set, where $M$ is real
symmetric, $r$ is bounded, and $0<c\le p_i=\Im r_i\le C$.
Assume $\ker D=\CC\mathbf1$. Suppose $q_i=\Im g_{ii}\le C$, the lag
$e'=\diag(g)-r$ satisfies $\|e'\|_\infty=o(t^{.36\delta})$ uniformly in
the construction parameters, and
\[
 \|K\|<\tfrac12,\qquad
 \|(\Id-K)^{-1}\|_{\infty\to\infty}\le C\log(1/t).
\]
For small $t$ there is a measurable vector $\zeta$, with
$\|\zeta\|_\infty\le t^{.35\delta}$, such that for $r'=r+\zeta$ and
$g'=(M-t\diag r')^{-1}$ one has $g'_{ii}=r'_i$.
Writing $\rho=\|\zeta\|_\infty$, the changes satisfy
\begin{align}
 \max_i\|\Delta h_i\|_2&\le C\rho,
 &\max_{i,j}|\Delta h_{ij}|&\le C\sqrt t\,\rho,
 &\max_i\|\Delta h_i\|_4&\le Ct^{1/4}\rho,
 \label{eq:assembly-reset-rows}\\
 \|K'-K\|&\le C\rho,
 &(1-C\rho)D&\le D'\le(1+C\rho)D.
 \label{eq:assembly-reset-forms}
\end{align}
Here $D'$ uses $p'=\Im r'$, which remains in a fixed positive bounded
range. The corresponding inverse forms on neutral charges change by
a factor $1+O(\rho)$.
\end{proposition}

\begin{proof}
The dissipative inverse bound gives $\|g\|\le C/t$, and Ward gives
\[
 \|g_{i\cdot}\|_2^2
 =t^{-1}\sum_j|h_{ij}|^2\le C/t.
\]
For a vector $z$ in a sufficiently small fixed $\ell^\infty$ ball,
put $Z=\diag z$ and $g_z=(M-t\diag(r+z))^{-1}$.
The Neumann series converges because $\|tZg\|\le C\|z\|_\infty$.
Its first diagonal derivative is
$t\sum_jg_{ij}^2z_j=(Kz)_i$. Thus
\begin{equation}
 \diag(g_z)=\diag(g)+Kz+\mathcal R(z),\qquad
 \|\mathcal R(z)\|_\infty\le C\|z\|_\infty^2.
 \label{eq:assembly-reset-taylor}
\end{equation}
For completeness, every diagonal term of order $k\ge2$ is bounded
using the two endpoint rows and $k-1$ intervening operator norms:
\[
 t^k(Ct^{-1/2})^2(Ct^{-1})^{k-1}\|z\|_\infty^k
 \le C^{k+1}\|z\|_\infty^k.
\]
The same expansion, replacing one factor of $z$ by $z-w$, proves
\[
 \|\mathcal R(z)-\mathcal R(w)\|_\infty
 \le C a\|z-w\|_\infty
 \quad\text{if }\|z\|_\infty,\|w\|_\infty\le a.
\]
There is no volume or negative power of $t$ in either constant.

Set $P=(\Id-K)^{-1}$. The reset equation is
\[
 z=P(e'+\mathcal R(z)).
\]
Take $a=t^{.35\delta}$. The image of the radius-$a$ ball has radius
at most
\[
 C\log(1/t)\bigl(o(t^{.36\delta})+Ct^{.70\delta}\bigr)=o(a),
\]
and its Lipschitz constant is at most
$C\log(1/t)t^{.35\delta}=o(1)$.
It is therefore a contraction. Its unique fixed point in that ball is
$\zeta$. Iteration from zero proves measurability in the initial data.
For small $t$, its imaginary part cannot move $p$ out of a fixed
positive bounded range.

The exact resolvent identity, in the $h$ normalization, is
\[
 \Delta h=h\sqrt t\,\diag(\zeta)h'.
\]
Both old and new inverse operator norms are $O(t^{-1})$; their endpoint
row norms in the $h$ normalization are bounded. For the new rows this
follows also by multiplying an old row by the convergent Neumann factor.
The identity gives $\|\Delta h_i\|_2\le C\rho$, and Cauchy--Schwarz
using its two endpoint rows gives $|\Delta h_{ij}|\le C\sqrt t\,\rho$.
Interpolating these two estimates proves
$\|\Delta h_i\|_4\le Ct^{1/4}\rho$.
Since
\[
 \sum_j|K'_{ij}-K_{ij}|
 \le \|\Delta h_i\|_2(\|h'_i\|_2+\|h_i\|_2)\le C\rho,
\]
symmetry and the row-sum bound prove the norm assertion for $K$.

For a real diagonal potential $f$, put $B=\sqrt t\,\diag(\zeta)$.
Then $h'=(h^{-1}-B)^{-1}$ and $[f,B]=0$, so
$[f,h']=h'h^{-1}[f,h]h^{-1}h'$.
Factoring the outside matrices gives the exact identity
\begin{equation}
 [f,h']
 =(\Id-h\sqrt t\,\diag\zeta)^{-1}
    [f,h]
   (\Id-\sqrt t\,\diag\zeta\,h)^{-1}.
 \label{eq:assembly-reset-commutator}
\end{equation}
Write the two outside factors as $L$ and $R$. Each, and each inverse,
differs from identity by $O(\rho)$ in operator norm, since
$\|h\|=O(t^{-1/2})$. Keep the nonconstant $p$ factors explicitly:
if $P_p=\diag(\sqrt p)$ and $C_f=P_p[f,h]P_p$, then
\[
 P_p[f,h']P_p
   =(P_pLP_p^{-1})C_f(P_p^{-1}RP_p).
\]
The outside factors and their inverses are $\Id+O(\rho)$, because
$P_p$ and $P_p^{-1}$ have bounded norm. Their multiplication changes
the Hilbert--Schmidt norm by two-sided factors $1+O(\rho)$.
Finally $p'_i/p_i=1+O(\rho)$ uniformly, so replacing $P_p$ by
$P_{p'}$ has the same relative cost. Because
\[
 2D(f)=\|\sqrt p\,[f,h]\sqrt p\|_{\mathrm{HS}}^2,
\]
we obtain \eqref{eq:assembly-reset-forms}.
These two-sided inequalities preserve the constant kernel, and inversion
on the neutral subspace proves the final assertion.
\end{proof}

\subsection{Closure of the deterministic induction}

We can now return from a completed stage to the same assumptions with
which the next stage begins. The probabilistic work in the next section
will control how many pending sites remain; the result here applies
pathwise to every finite history produced by the trial rules, with
numerical rollback on failure.

\begin{proposition}[Uniform stage bounds]
\label{prop:deterministic-induction}
Choose the fixed energy window and the initialization constants as in
Proposition~\ref{prop:initialization}. For sufficiently large $n_0$,
uniformly in $m\ge4n_0$ and in the fixed energy in that window, the
following statements hold along every finite history produced by the
ordinary-sweep and clearing rules, with failed trials rolled back.
At each completed stage start the diagonal fixed-point relation,
the three comparisons \eqref{eq:comparators}, the fourth-row bound
$\max_i\sum_j|h_{ij}|^4\le t^{1+\delta}$, and $\|K\|\le1/3$ hold.
The vectors $r$ remain bounded with their imaginary parts in one fixed
positive bounded range. The low comparison constants remain in fixed
positive bounded ranges, and the rough high constant $C_0$ can be
chosen independently of the stage. The kernel of $D$ consists of the
constants, at least $N/2$ sites remain active, and the terminal metric
satisfies \eqref{eq:assembly-terminal-metric}.

During a clearing stage, every committed prestate has lag
$o(t^{.36\delta})$, fourth-row sums below $.8t^{1+\delta}$, and the
comparison bounds furnished by Proposition~\ref{prop:committed-assembly}.
No comparison for $D$ or bound on $K$ is asserted at a provisional heat
prefix. All constants and smallness thresholds are independent of the
history and of the torus volume.
\end{proposition}

\begin{proof}
Initialization supplies the assertions at the first stage, with
$\|K\|=o(1)$. Enlarge the permitted bounds on $r$, $p$, and the low
comparison constants by fixed factors, retaining positive lower bounds.
We show that the estimates obtained under these enlarged ranges remain
inside them. The form equivalence $D\asymp W$ used below follows from
the upper rough comparison for $X$, and from the two lower bounds
$X\ge cb\Lambda^+$ and $X\ge c\Id$.

Apply Proposition~\ref{prop:ordinary-sweep} at the old scale $t_{\rm old}$,
and write $t=\theta t_{\rm old}$. On neutral physical-site charges define
the comparison form
\begin{equation}
 \mathcal X=\theta X_{\rm old}
  +(1-\theta)\diag(p^{-2}),
 \label{eq:assembly-affine-reset}
\end{equation}
The exact reserve-port identity in Lemma~\ref{lem:duplication}
identifies $\mathcal X$ with $\theta X_0$ restricted to reserve
charges, where $X_0$ is the duplicated stage-start inverse form.
Proposition~\ref{prop:ordinary-sweep} places the actual post-sweep
inverse form between $(1-\epsilon)\mathcal X$ and
$(1+\epsilon)\mathcal X$, with $\epsilon=Ct_{\rm old}^{.10\delta}$.
The low scale changes from $b_{\rm old}$ to $\theta b_{\rm old}$.
The physical lag after the sweep is $O(t^{.42\delta})$, its fourth-row
sums are below $.7t^{1+\delta}$, and its $K$ norm has increased by at
most $Ct_{\rm old}^{.10\delta}$.

If this stage is within the clearing cutoff, consider its snapshot
components in their prescribed order. The post-sweep state satisfies
\eqref{eq:assembly-initial-data}: its lag and row margins were just
obtained, and the form equivalence supplies its comparison. At the first
trial, and then at any trial whose preceding prefix has the bounds of
Proposition~\ref{prop:committed-assembly}, the lag and row hypotheses of
Proposition~\ref{prop:clearing-trial} hold. The location induction of
Lemma~\ref{lem:clearing-locality} applies throughout the stopped trial.
On DONE, the heat and core calculation following
Proposition~\ref{prop:committed-assembly} supplies its two operation inputs
for the new numerical endpoint. Applying that proposition to the enlarged
prefix gives a total relative form error $O(t^{.10\delta})$.
Its proof bounds the total $K$ increment by
$Ct^{.125\delta-d\nu}$ and the final lag by
$Ct^{.42\delta-d\nu}=o(t^{.36\delta})$.
In fourth-row norm the cumulative change, divided by
$t^{(1+\delta)/4}$, is $O(t^{.0225\delta-d\nu})=o(1)$.
For large $n_0$ it moves the sweep's $.7$ bound to a bound strictly
below $.8$, uniformly at every committed prefix.
These bounds supply the starting margins for the next component.
Failed trials restore the previous analytic state; their new pending
marks do not alter any of these bounds. If no clearing is scheduled,
the sweep endpoint itself has the stronger required estimates.

The committed endpoint has $\|K\|<1/2$ once the total summable increments
are chosen small. Lemma~\ref{lem:inverse-row}, applied with its lag and
the terminal metric, gives
$\|(\Id-K)^{-1}\|_{\infty\to\infty}\le C\log(1/t)$.
Proposition~\ref{prop:reset} now restores $g_{ii}=r_i$.
Its fourth-row norm cost relative to $t^{(1+\delta)/4}$ is
\[
 Ct^{1/4}t^{.35\delta}/t^{(1+\delta)/4}
 =Ct^{.10\delta}=o(1).
\]
The full fourth-row bound is therefore restored. Its form error is
$O(t^{.35\delta})$, and the change of $r$ and the additional $K$ cost
are bounded by the same power.

It remains to check the three individual comparison constants.
For the sharp upper bound, \eqref{eq:assembly-affine-reset} transforms
the diagonal coefficient $1+t_{\rm old}^{\kappa}$ into
$1+\theta t_{\rm old}^{\kappa}$. Clearing, sweep, and reset change this
by $O(t_{\rm old}^{.10\delta})$; the change from $p$ to $p'$ at reset
has the smaller error $O(t^{.35\delta})$.
For $\theta\in[.25,.625]$, continuity and $\kappa<1$ give
\[
 \inf_{\theta\in[.25,.625]}(\theta^\kappa-\theta)>0.
\]
Thus, for small $t_{\rm old}$,
\[
 1+\theta t_{\rm old}^{\kappa}
       +Ct_{\rm old}^{.10\delta}
 \le 1+(\theta t_{\rm old})^\kappa,
\]
because $.10\delta>\kappa=.005\delta$.
For the rough upper bound the high coefficient is at most
\[
 (1+o(1))\bigl(\theta C_0+C\bigr).
\]
Choosing $C_0$ large first and $n_0$ large second keeps this below $C_0$,
since $\theta\le.625$. Both low comparison coefficients change only by
factors $1+O(t_{\rm old}^{.10\delta})$; the lower factor is interpreted
as $1-O(t_{\rm old}^{.10\delta})$.

All these errors are bounded by a constant times $4^{-c n}$ for a
fixed $c>0$. Their sum from $n_0$ onward can be made arbitrarily small,
and the corresponding positive products stay bounded above and away
from zero. This closes the enlarged ranges for the low constants,
$r$, and $p$. It also keeps $\|K\|\le1/3$, using the initialization
margin and the summable norm increments. Bounded $p,q$ give
$\|D\|\le C$, hence $X\ge c\Id$ on neutral charges. Relative form
comparisons and Schur complementation preserve the one-dimensional
constant kernel. The geometry and active-count bounds are supplied
independently by Lemma~\ref{lem:clearing-locality}.
After the clearing cutoff the terminal set no longer changes, and
the permissible metric defect $t^{-1.03}$ only increases as $t$
decreases, so the same metric bound continues to hold.

Finally, if $a_\lambda=(\lambda^2/3)4^{n_0}$, then
$a_\lambda\in(1/4,1]$ and
$t_n=a_\lambda(4^{-n}+4^{-m})$.
The geometric sums and smallness choices above are therefore uniform in
$m$ and in the energy after the initialization window has been fixed.
One fixed dimension-dependent lower threshold for $n_0$ suffices
simultaneously for all $m\ge4n_0$, $a_\lambda\in(1/4,1]$, and $E\in I_d$.
For $d\ge4$, the intrinsic choice of the final disorder threshold will be made in
Definition~\ref{def:threshold} after the probability estimates.
The argument proves preservation along every finite history produced by the rules;
almost-sure termination of ordinary retries is a separate consequence
of their conditional failure bound, used in the next section.
\end{proof}
